% Required preamble commands/packages:
% \usepackage{amsmath,amssymb,amsthm}
% \newtheorem{theorem}{Theorem}
% \newtheorem{corollary}{Corollary}
% \newcommand{\E}{\mathbb{E}}
% \newcommand{\R}{\mathbb{R}}
% \newcommand{\Tr}{\operatorname{tr}}
% \newcommand{\diag}{\operatorname{diag}}
% \newcommand{\argmin}{\operatorname*{arg\,min}}

\section{Theoretical Foundation}

\subsection{Problem Setup}

We consider an idealized linear-Gaussian state-space model
\begin{align}
    x_{t+1} &= A x_t + w_t, \qquad w_t \sim \mathcal{N}(0, Q), \\
    y_t &= C x_t + v_t, \qquad v_t \sim \mathcal{N}(0, R),
\end{align}
where $x_t \in \R^d$ is the latent state, $y_t$ is the observation, and $A$
is a stable transition matrix. The model is meant as a tractable analogue of
the deterministic component of an RSSM transition, in which the recurrent state
evolves under a learned dynamics operator with both slowly and quickly varying
modes.

We assume that the transition has a slow-fast spectral structure. Let
$\lambda_i(A)$ denote the eigenvalues of $A$. For a threshold
$\lambda_{\mathrm{gap}}$, define the true slow subspace as
\begin{equation}
    \mathcal{S}_{\mathrm{true}}
    =
    \operatorname{span}\{q_i : |\lambda_i(A)| > \lambda_{\mathrm{gap}}\},
\end{equation}
where $q_i$ is the eigenvector or invariant-subspace direction associated with
$\lambda_i(A)$. Intuitively, modes with $|\lambda_i(A)|$ close to one persist
over time and therefore generate low temporal derivative energy.

The corresponding slow-feature objective is
\begin{equation}
    \min_{U \in \R^{d \times r}}
    \E\left[\left\|\frac{d}{dt} U^\top x_t\right\|_2^2\right]
    \quad
    \text{subject to}
    \quad
    \E\left[U^\top x_t x_t^\top U\right] = I_r.
    \label{eq:slow_feature_objective}
\end{equation}
In discrete time we use $\dot{x}_t \approx x_{t+1} - x_t$, so the objective
penalizes the derivative variance of the projected latent trajectory while the
constraint prevents the degenerate zero solution.

\subsection{Theorem 1 --- Slow Subspace Identifiability}

\begin{theorem}[Slow subspace identifiability]
\label{thm:slow_subspace_identifiability}
Consider the linear-Gaussian model above. Assume:
\begin{enumerate}
    \item The latent dynamics are linear and stable: $x_{t+1}=Ax_t+w_t$ with
    $\rho(A)<1$.
    \item The process noise is Gaussian with nonsingular covariance, and the
    stationary covariance $\Sigma_x=\E[x_t x_t^\top]$ exists and is positive
    definite.
    \item The transition admits a slow-fast decomposition with a nonzero
    spectral gap: for the $r$ slow modes and the remaining fast modes,
    $|1-\lambda_i|^2 < |1-\lambda_j|^2$ whenever $i$ is slow and $j$ is fast.
    \item Sufficient data are available so that the empirical covariance and
    derivative covariance converge to their population values.
\end{enumerate}
Then any global minimizer $U_\star$ of the constrained slow-feature objective
in~\eqref{eq:slow_feature_objective} spans the true slow subspace:
\begin{equation}
    \operatorname{span}(U_\star) = \mathcal{S}_{\mathrm{true}},
\end{equation}
up to an invertible rotation within the $r$-dimensional slow subspace.
\end{theorem}

\paragraph{Proof sketch.}
Write the discrete derivative as $\dot{x}_t=x_{t+1}-x_t=(A-I)x_t+w_t$.
The slow-feature objective is therefore a Rayleigh quotient involving the
derivative covariance
\begin{equation}
    \Sigma_{\dot{x}} = \E[\dot{x}_t \dot{x}_t^\top],
\end{equation}
under the whitening constraint induced by $\Sigma_x$. After whitening the
stationary covariance, the problem reduces to selecting the eigenvectors of the
whitened derivative covariance with the smallest eigenvalues.

Under the slow-fast spectral gap assumption, the derivative variance associated
with a transition eigenmode scales with $|1-\lambda_i|^2$ plus the contribution
of process noise. Modes whose eigenvalues are closest to one therefore have the
smallest derivative variance. The gap ensures that the $r$ smallest generalized
eigenvalues are exactly the slow modes, so the minimizer recovers their span.

Finally, if multiple slow modes have equal derivative variance, the objective
identifies their shared subspace rather than a unique basis. This is why the
recovery statement is up to rotation within $\mathcal{S}_{\mathrm{true}}$.
TODO: In the final proof, state the exact diagonalizability/normality and
noise-isotropy conditions needed for this eigenmode argument to hold without
cross-mode terms.

\paragraph{Remark: empirical PCA plus autocorrelation.}
The empirical procedure used in our diagnostics is a finite-sample proxy for
the theorem. PCA first restricts attention to high-variance latent directions,
and lag-1 autocorrelation then ranks temporal projections by persistence. In
the idealized linear-Gaussian setting, high persistence corresponds to small
derivative variance. In trained nonlinear RSSMs this equivalence need not hold
exactly, which motivates the empirical alignment checks in Section~4.

\subsection{Damping Corollary}

\begin{corollary}[Slow-subspace damping]
\label{cor:damping}
Assume the conditions of Theorem~\ref{thm:slow_subspace_identifiability}, and
suppose $U$ is an orthonormal basis aligned with the true slow drift subspace.
Let $\lambda_s^{\mathrm{true}}$ denote the relevant slow eigenvalue of the true
latent dynamics and $\lambda_s^{\mathrm{model}}$ the corresponding slow
eigenvalue of the learned imagination dynamics. Applying damping with rate
$\eta$ along $U$ reduces slow-subspace imagination MSE if and only if
\begin{equation}
    0 < \eta
    <
    2\left(1 -
    \frac{\lambda_s^{\mathrm{true}}}{\lambda_s^{\mathrm{model}}}
    \right),
    \label{eq:damping_bound}
\end{equation}
assuming $\lambda_s^{\mathrm{model}}>\lambda_s^{\mathrm{true}}$ and both slow
eigenvalues are real and positive.
\end{corollary}

\paragraph{Proof sketch.}
In the aligned slow basis, damping scales the learned slow transition by a
factor of $(1-\eta)$. The damped model eigenvalue is therefore
$(1-\eta)\lambda_s^{\mathrm{model}}$. Damping improves the slow-subspace
prediction error precisely when this damped eigenvalue is closer to the true
slow eigenvalue than the undamped model eigenvalue:
\begin{equation}
    \left|(1-\eta)\lambda_s^{\mathrm{model}}
    - \lambda_s^{\mathrm{true}}\right|
    <
    \left|\lambda_s^{\mathrm{model}}
    - \lambda_s^{\mathrm{true}}\right|.
\end{equation}
Solving this scalar inequality yields the upper bound in
\eqref{eq:damping_bound}. The upper limit prevents over-correction: too much
damping flips the slow update past the true transition and increases error.

\paragraph{Remark: alignment caveat.}
This result assumes $U$ aligns with the true drift subspace. Our empirical
findings (Section~4) reveal that posterior-estimated bases may not satisfy this
condition, motivating the use of empirically estimated drift bases.

\subsection{Limitations of the Linear-Gaussian Framework}

The theory does not cover the full complexity of modern RSSMs. It omits
nonlinear transition functions, tanh saturation, learned stochastic latent
interactions, action-conditioned dynamics, nonstationary rollout distributions,
and decoder-induced distortions. It also treats the relevant subspaces as
linear and time-invariant, whereas trained world models may rotate or reshape
their drift directions across training.

The purpose of the theory is therefore not to prove that SMAD must work in a
trained Dreamer model. Instead, it identifies a clean regime in which slow
features are recoverable and slow-subspace damping is beneficial. The following
sections test whether the qualitative predictions of this framework hold in
nonlinear trained models.
